% Authored. Numbers come from generated/macros.tex.

\begin{abstract}
\noindent
A general-government balance is reported as a single number, although it is the sum
of institutionally distinct subsectors with different revenue bases, expenditure
mandates and borrowing constraints. We assemble a harmonised annual panel of net
lending and net borrowing for Portuguese General Government and its three
subsectors over \PanelYears\ years, \PanelStart--\PanelEnd, and decompose the
aggregate exactly into its subsector components and their annual changes.

\noindent
The composition is asymmetric and stable. The aggregate balance is positive in
\NumPositiveYears\ years (\PositiveYearList); the Central Government balance is
negative in all \CentralNegativeYears; the Social Security balance is positive in
\SsfPositiveYears, with a longest run of \SsfLongestPositiveRun\ years. General
Government and Central Government track each other closely, at a correlation of
\AggregateCentralCorrelation\ and a median absolute gap of
\AggregateCentralMedianGap\ percentage points of GDP. The empirically informative
feature of the \NumPositiveYears\ surplus years is that the combined
non-Social-Security balance is negative in all of them; that the Social Security
balance then exceeds it in absolute size follows from the identity.

\noindent
Two findings cut against readings the aggregate invites. First, a persistently
negative Central Government B.9 does not imply a persistently negative primary
balance: excluding interest, the primary balance is positive in
\PrimaryPositiveYears\ of the \PrimaryObservedYears\ years with detailed accounts,
from \PrimaryPositiveFirst\ to \PrimaryPositiveLast. Second, the ESA 2010 Social
Security Funds sector and the budget-execution systems published by the Portuguese
Public Finance Council are different accounting objects whose balances differ in
every one of the \SsfBoundaryYears\ years they overlap, so quoting one for the
other misstates the figure that enters the national accounts.

\noindent
Applying the same definitions to \BenchmarkCountries\ European reporters shows the
composition is distinctive in a specific respect. A permanently deficit-running central
tier is common --- \CentralAlwaysNegativeCountries\ reporters record one in every year
--- but Portugal's Social Security surplus is in the upper tail on both frequency and
size, and it is the only reporter whose every aggregate-surplus year combines that
surplus with a negative non-Social-Security balance, a composition found in
\SurplusOffsettingShare\% of European surplus years.

\noindent
Magnitudes are regime-specific: the aggregate averages \GgMeanHistorical\% of GDP
before the 1995 statistical splice and \GgMeanModern\% after, so no level statistic
is pooled across it. The analysis is descriptive and accounting-based; it decomposes
recorded quantities and identifies no causal effect.

\medskip
\noindent
\textbf{Keywords:} general government; institutional subsectors; social security
funds; accounting decomposition; fiscal statistics; Portugal

\noindent
\textbf{JEL:} H61, H62, H55, E62
\end{abstract}
